\documentclass[10pt,letterpaper]{article}
\usepackage[utf8]{inputenc}
\usepackage[T1]{fontenc}
\usepackage[margin=0.65in]{geometry}
\usepackage{amsmath,amssymb}
\usepackage{booktabs}
\usepackage{longtable}
\usepackage{graphicx}
\usepackage{hyperref}
\hypersetup{colorlinks=true,linkcolor=blue,citecolor=blue,urlcolor=blue}

\title{\textbf{Beetle Tree of Life: Permian--Cenozoic Coleoptera Phylogenomics (308.0 to 0.0 Ma; Crown Coleoptera 292.0 Ma), Suborder Divergence, and Microbial PCWDE Horizontal Gene Transfer (GH28, GH45, GH5)}}
\author{\textbf{Bryan P. White} \\ CRV Science Labs \\ \texttt{bryan@crvscience.com}}
\date{November 7, 2026 \quad $\cdot$ \quad Preprint ID: PREPRINT-13 (v2.0.0 Empirical Preprint) \quad $\cdot$ \quad DOI: 10.5281/zenodo.24031323}

\begin{document}
\maketitle

\begin{abstract}
\noindent \textbf{Abstract ---} Analyzing 101 verified NCBI RefSeq and GenBank sequence records totaling 1,378,439 bp on disk---comprising 83 complete Coleoptera, Strepsiptera, and Neuropterida mitochondrial genomes (1,354,809 bp across circular phylogram tips T001--T083; aligned across 35,624 full-length columns by *MAFFT v7.526 in 896.99 s, 197,673 full-length columns by Native Rust gotoh-msa in 1.00 s, and 4,125 bp of conserved COX1, COB, COX2, and COX3 codons) plus 18 nuclear glycoside hydrolase GH28 polygalacturonase, GH45 endoglucanase, and GH5 mannanase mRNA transcripts (23,630 bp across NUC084--NUC101; 3,321 full-length *MAFFT columns in 2.08 s)---across a 32-chain full-length combinatorial model sweep and 1,000 *IQ-TREE v2.4.0 ultrafast bootstrap replicates~\cite{White2026RuNIE,McKenna2019,Kirsch2014}, we date the Neuropterida + Coleopterida root (Node N81) to 308.00 Ma [95\% CI: 308.00--328.00 Ma] (100\% UFBoot) and Crown Coleoptera (Node N77) to 292.00 Ma [292.00--305.00 Ma] (100\% UFBoot), resolving Myxophaga + Adephaga (Node N75: 258.00 Ma, 98\% UFBoot) sister to Archostemata + Polyphaga (Node N62: 286.00 Ma, 89\% UFBoot). Felsenstein ModelTest auditing shows that 3rd-codon RY3 recoding and Dayhoff-6 amino-acid recoding reduce Relative Composition Frequency Variability from 0.0624 (whole mitogenome) and 0.0312 (4-gene 3rd codons) down to 0.0074 (-76.3\%) and 0.0034 (-89.1\%), improving marginal log-likelihood by +42,036.71 nats on the 5,727 bp total-evidence matrix ($\ln L = -140,326.15$ vs.\ $-182,362.86$) and +93,310.59 nats on the 35,624-column full-length mitogenome alignment ($\ln L = -656,679.01$ vs.\ $-749,989.60$). Dual-Pass SHA-256 blinded K2P permutation auditing ($B = 500$ across 3,403 pairs) yields $Z_{\text{blind}} = 3.81$ and $D_{\text{ctx}} = 0.0000$ while autonomously detecting the 100.00\% identical RefSeq/GenBank duplicate pair T007 (\texttt{NC\_011328.1}) and T071 (\texttt{EU877953.1}) in \textit{Tetraphalerus bruchi}. Finally, *PAML v4.10.10 \texttt{codeml} selection analysis across GH28 (350 codons, $\ln L_0 = -4,701.50$, $\omega_0 = 0.0484$), GH45 (232 codons, $\ln L_0 = -3,025.53$, $\omega_0 = 0.2176$; M2 $2\Delta\ln L = 14.729$, $p = 1.24 \times 10^{-4}$), and GH5 (385 codons, $\ln L_0 = -2,968.40$, $\omega_0 = 0.0852$) confirms strong purifying selection preserving the microbial catalytic triads following Early Jurassic acquisition in Crown Phytophaga (Node N47: 182.00 Ma [164.08--182.00 Ma], 98\% UFBoot).
\end{abstract}

\section{Dataset Architecture \& Multi-Regime Full-Length Alignment Telemetry}
Across all 101 verified NCBI RefSeq and GenBank records (1,378,439 bp on disk: 83 complete mitogenomes = 1,354,809 bp across T001--T083 and 18 nuclear PCWDE mRNAs = 23,630 bp across NUC084--NUC101)~\cite{White2026RuNIE,McKenna2019,Kirsch2014,Zhang2018}, external *MAFFT v7.526 aligned the 83 full-length mitogenomes into 35,624 columns in 896.99 s and the 18 nuclear transcripts into 3,321 columns in 2.08 s (38,945 combined Total-Evidence columns in 899.07 s)~\cite{Katoh2013}. In 100\% Native Rust (Regimes 1A and 1B), \texttt{gotoh-msa} aligned the 83 full-length mitogenomes into 197,673 columns in 1.00 s (RY3) and 1.17 s (Raw 4-state). Extracting the 4 conserved mitochondrial protein-coding genes (COX1 = 514 codons, COB = 375 codons, COX2 = 227 codons, COX3 = 259 codons) yielded a 4,125 bp (1,375-codon) matrix (Pipeline B) and a 5,727 bp Total-Evidence matrix (Pipeline A).

\section{Permian--Cenozoic LSD2 Chronogram \& Compositional RCFV Rescue}
Across the 83 mitogenomes, 3rd-codon GC content spans from 11.43\% in \textit{Trachypachus holmbergi} (\texttt{NC\_011329.1}, GC = 20.54\%) and 15.51\% in \textit{Xenos peckii} (\texttt{CM079275.1}, GC = 14.55\%) up to 33.03\% in \textit{Tetraphalerus bruchi} (\texttt{NC\_011328.1} / \texttt{EU877953.1}, GC = 33.01\%) and 32.82\% in \textit{Lucanus mazama} (\texttt{NC\_013578.1}, GC = 32.87\%)~\cite{Niehuis2012,Sheffield2008,Sheffield2009}. Stage 3 RY3 recoding reduces RCFV from 0.0624 (whole mitogenome) and 0.0312 (4-gene Pos3) to 0.0074 (-76.3\%), and Dayhoff-6 recoding reduces RCFV to 0.0034 (-89.1\%)~\cite{Susko2007}. Calibrated *IQ-TREE v2.4.0 \texttt{--date} LSD2 dating~\cite{Minh2020,To2016} places the Neuropterida+Coleopterida root (N81) at 308.00 Ma [308.00--328.00 Ma], Crown Coleoptera (N77) at 292.00 Ma [292.00--305.00 Ma], Myxophaga+Adephaga (N75) at 258.00 Ma [258.00--273.77 Ma], Crown Adephaga (N74) at 235.00 Ma, Crown Myxophaga (N64) at 212.37 Ma, Archostemata+Polyphaga (N62) at 286.00 Ma, Crown Archostemata (N61) at 222.46 Ma, Core Polyphaga (N58) at 246.04 Ma, Crown Cucujiformia (N57) at 210.00 Ma, Crown Phytophaga (N47) at 182.00 Ma [164.08--182.00 Ma], Crown Curculionoidea (N46) at 162.00 Ma, and Crown Chrysomeloidea (N33) at 156.00 Ma~\cite{Kirejtshuk2014,Cai2022,Shin2018}.

\section{Microbial-to-Beetle PCWDE Horizontal Gene Transfer \& *PAML v4.10.10 Selection}
Whereas \textit{Tribolium castaneum} (\texttt{NC\_003081.2}) encodes only ancestral metazoan GH9 endoglucanase~\cite{Tribolium2008}, Crown Phytophaga (N47: 182.00 Ma) acquired microbial GH28 polygalacturonases (\texttt{XM\_076418249.1}--\texttt{XM\_076416151.1}), GH45 inverting cellulases (\texttt{XM\_076407000.1}, \texttt{XM\_023173738.2}--\texttt{XM\_023173742.2}), and GH5 mannanases (\texttt{XM\_076408428.1}, \texttt{XM\_076400312.1}--\texttt{XM\_076400308.1})~\cite{Kirsch2014,Pauchet2010,McKenna2016,Acuna2012,Schoville2018,Keeling2013,Parisot2021,Hazzouri2020}. Executing *PAML v4.10.10 \texttt{codeml}~\cite{Yang2005} confirms strong purifying selection across GH28 ($N = 5$, 350 codons, $\ln L_0 = -4,701.50$, $\omega_0 = 0.0484$, $\kappa = 1.441$; conserving NTDGFD, QDDCVA, GGHGLS, and NGIH[I/V]KTH in template \texttt{1BHE}~\cite{vanSanten1999,Zhang2022USalign}), significant clade-specific divergence in GH45 ($N = 7$, 232 codons, $\ln L_0 = -3,025.53$, $\omega_0 = 0.2176$; M2 $\ln L_2 = -3,018.17$, $\omega_{\text{Curc}} = 0.0056$ vs.\ $\omega_{\text{Chrys}} = 0.2651$, $2\Delta\ln L = 14.729$, $p = 1.24 \times 10^{-4}$), and strong constraint in GH5 ($N = 6$, 385 codons, $\ln L_0 = -2,968.40$, $\omega_0 = 0.0852$, $\kappa = 1.077$)~\cite{Kirsch2016,Morcos2011}.

\begin{thebibliography}{36}
\bibitem{White2026RuNIE} White, B. P. (2026). RuNIE (Rust Native Inference Engine): An Open-Source Architecture for Autonomous Bioinformatics Pipeline Synthesis, Blinded Validation, and Bayesian Scope-Depth Auditing. \textit{CRV Science Labs Preprint}, v2.0.1. \url{https://doi.org/10.5281/zenodo.23969275}
\bibitem{McKenna2019} McKenna, D. D., et al. (2019). The evolution and genomic basis of beetle diversity. \textit{PNAS}, 116(49), 24729--24737. \url{https://doi.org/10.1073/pnas.1909655116}
\bibitem{Kirsch2014} Kirsch, R., et al. (2014). Horizontal gene transfer facilitated the evolution of pectinolytic ability and herbivory in beetles. \textit{Nature}, 513, 377--382. \url{https://doi.org/10.1038/nature13681}
\bibitem{Niehuis2012} Niehuis, O., et al. (2012). Genomic and morphological evidence converge to resolve the enigma of Strepsiptera. \textit{Current Biology}, 22(14), 1309--1313. \url{https://doi.org/10.1016/j.cub.2012.05.018}
\bibitem{Sheffield2008} Sheffield, N. C., et al. (2008). A comparative analysis of mitochondrial genomes in Coleoptera. \textit{Mol. Biol. Evol.}, 25(11), 2499--2509. \url{https://doi.org/10.1093/molbev/msn198}
\bibitem{Sheffield2009} Sheffield, N. C., et al. (2009). Nonstationary evolution and compositional heterogeneity in beetle mitochondrial genomes. \textit{Syst. Biol.}, 58(4), 381--394. \url{https://doi.org/10.1093/sysbio/syp038}
\bibitem{Zhang2018} Zhang, S. Q., et al. (2018). Evolutionary history of Coleoptera revealed by extensive sampling of genes and species. \textit{Nature Communications}, 9, 205. \url{https://doi.org/10.1038/s41467-017-02644-4}
\bibitem{Misof2014} Misof, B., et al. (2014). Phylogenomics resolves the timing and pattern of insect evolution. \textit{Science}, 346(6210), 763--767. \url{https://doi.org/10.1126/science.1257570}
\bibitem{Tribolium2008} Tribolium Genome Sequencing Consortium (2008). The genome of the model beetle and pest \textit{Tribolium castaneum}. \textit{Nature}, 452, 949--955. \url{https://doi.org/10.1038/nature06784}
\bibitem{Kirejtshuk2014} Kirejtshuk, A. G., et al. (2014). Evolution of the elytral venation and structural adaptations in the oldest Palaeozoic beetles. \textit{J. Syst. Palaeontol.}, 12(5), 575--600. \url{https://doi.org/10.1080/14772019.2013.821530}
\bibitem{Cai2022} Cai, C., et al. (2022). Integrated phylogenomics and fossil data illuminate the evolution of beetles. \textit{R. Soc. Open Sci.}, 9(3), 211771. \url{https://doi.org/10.1098/rsos.211771}
\bibitem{Pauchet2010} Pauchet, Y., et al. (2010). Diversity of beetle genes encoding novel plant cell wall degrading enzymes. \textit{Insect Biochem. Mol. Biol.}, 40(1), 60--71. \url{https://doi.org/10.1016/j.ibmb.2009.12.007}
\bibitem{Shin2018} Shin, S., et al. (2018). Phylogenomic data yield new and robust insights into the phylogeny and evolutionary history of weevils. \textit{Mol. Biol. Evol.}, 35(4), 823--836. \url{https://doi.org/10.1093/molbev/msx324}
\bibitem{McKenna2016} McKenna, D. D., et al. (2016). Genome of the Asian longhorned beetle (\textit{Anoplophora glabripennis}). \textit{Genome Biology}, 17, 227. \url{https://doi.org/10.1186/s13059-016-1088-8}
\bibitem{Acuna2012} Acuna, R., et al. (2012). Adaptive horizontal transfer of a bacterial gene to an invasive insect pest of coffee. \textit{PNAS}, 109(11), 4197--4202. \url{https://doi.org/10.1073/pnas.1121190109}
\bibitem{vanSanten1999} van Santen, Y., et al. (1999). 1.68-A crystal structure of endopolygalacturonase II from \textit{Aspergillus niger}. \textit{J. Biol. Chem.}, 274(43), 30474--30480. \url{https://doi.org/10.1074/jbc.274.43.30474}
\bibitem{Kirsch2016} Kirsch, R., Heckel, D. G., \& Pauchet, Y. (2016). How the beetle became aba: The evolutionary fate of horizontally transferred pectinase genes in the Chrysomeloidea. \textit{BMC Evol. Biol.}, 16, 106. \url{https://doi.org/10.1186/s12862-016-0673-z}
\bibitem{Schoville2018} Schoville, S. D., et al. (2018). A model species for agricultural pest genomics: the genome of the Colorado potato beetle, \textit{Leptinotarsa decemlineata}. \textit{Sci. Rep.}, 8, 1931. \url{https://doi.org/10.1038/s41598-018-20154-1}
\bibitem{Keeling2013} Keeling, C. I., et al. (2013). Draft genome of the mountain pine beetle, \textit{Dendroctonus ponderosae}. \textit{Genome Biology}, 14, R27. \url{https://doi.org/10.1186/gb-2013-14-3-r27}
\bibitem{Parisot2021} Parisot, N., et al. (2021). The transposable element-rich genome of the cereal pest \textit{Sitophilus oryzae}. \textit{BMC Biology}, 19, 241. \url{https://doi.org/10.1186/s12915-021-01158-2}
\bibitem{Hazzouri2020} Hazzouri, K. M., et al. (2020). The genome of pest \textit{Rhynchophorus ferrugineus} reveals gene families important at the plant-beetle interface. \textit{Commun. Biol.}, 3, 323. \url{https://doi.org/10.1038/s42003-020-1060-8}
\bibitem{Hunt2007} Hunt, T., et al. (2007). A comprehensive phylogeny of beetles reveals the evolutionary origins of a superradiation. \textit{Science}, 318, 1913--1916. \url{https://doi.org/10.1126/science.1146954}
\bibitem{Farrell1998} Farrell, B. D. (1998). Inordinate fondness explained: why are there so many beetles? \textit{Science}, 281, 555--559. \url{https://doi.org/10.1126/science.281.5376.555}
\bibitem{Beutel2014} Beutel, R. G., et al. (2014). \textit{Insect Morphology and Phylogeny}. Walter de Gruyter. \url{https://doi.org/10.1515/9783110264043}
\bibitem{Katoh2013} Katoh, K., \& Standley, D. M. (2013). MAFFT multiple sequence alignment software version 7. \textit{Mol. Biol. Evol.}, 30(4), 772--780. \url{https://doi.org/10.1093/molbev/mst010}
\bibitem{Ranwez2018} Ranwez, V., et al. (2018). MACSE v2: toolkit for the alignment of coding sequences. \textit{Mol. Biol. Evol.}, 35(10), 2582--2584. \url{https://doi.org/10.1093/molbev/msy159}
\bibitem{Minh2020} Minh, B. Q., et al. (2020). IQ-TREE 2: new models and efficient methods for phylogenetic inference. \textit{Mol. Biol. Evol.}, 37(5), 1530--1534. \url{https://doi.org/10.1093/molbev/msaa015}
\bibitem{Lartillot2004} Lartillot, N., \& Philippe, H. (2004). A Bayesian mixture model for across-site heterogeneities in the amino-acid replacement process. \textit{Mol. Biol. Evol.}, 21(6), 1095--1109. \url{https://doi.org/10.1093/molbev/msh112}
\bibitem{Susko2007} Susko, E., \& Roger, A. J. (2007). On reduced amino acid alphabets for phylogenetic inference. \textit{Mol. Biol. Evol.}, 24(9), 2139--2150. \url{https://doi.org/10.1093/molbev/msm144}
\bibitem{Zhang2018ASTRAL} Zhang, C., et al. (2018). ASTRAL-III: polynomial time species tree reconstruction. \textit{BMC Bioinformatics}, 19, 153. \url{https://doi.org/10.1186/s12859-018-2129-y}
\bibitem{Yang2005} Yang, Z., Wong, W. S., \& Nielsen, R. (2005). Bayes empirical Bayes inference of amino acid sites under positive selection. \textit{Mol. Biol. Evol.}, 22(4), 1107--1118. \url{https://doi.org/10.1093/molbev/msi097}
\bibitem{Morcos2011} Morcos, F., et al. (2011). Direct-coupling analysis of residue coevolution captures native contacts. \textit{PNAS}, 108(49), E1293--E1301. \url{https://doi.org/10.1073/pnas.1111471108}
\bibitem{Zhang2022USalign} Zhang, C., et al. (2022). US-align: universal structure alignments of proteins, nucleic acids, and macromolecular complexes. \textit{Nature Methods}, 19(9), 1109--1115. \url{https://doi.org/10.1038/s41592-022-01585-1}
\bibitem{Heath2014} Heath, T. A., Huelsenbeck, J. P., \& Stadler, T. (2014). The fossilized birth-death process for coherent calibration of divergence-time estimates. \textit{PNAS}, 111(29), E2957--E2966. \url{https://doi.org/10.1073/pnas.1319091111}
\bibitem{To2016} To, T. H., et al. (2016). Fast dating using least-squares criteria and algorithms. \textit{Syst. Biol.}, 65(1), 82--97. \url{https://doi.org/10.1093/sysbio/syv068}
\bibitem{Eddy2011} Eddy, S. R. (2011). Accelerated profile HMM searches. \textit{PLOS Comput. Biol.}, 7(10), e1002195. \url{https://doi.org/10.1371/journal.pcbi.1002195}
\end{thebibliography}
\end{document}
